% !TeX root = anexo_semana02.tex
% !TeX program = pdflatex
% Anexo independiente de la semana 02: demostraciones directas.
\documentclass[10pt,a5paper]{article}
\newif\ifmostrarSoluciones
\mostrarSolucionestrue
\usepackage{estilo_apuntes}
\encabezado{Semana 2}{Anexo: demostraciones}

\begin{document}
\tituloclase{Primeros pasos en demostraciones}{Guía para demostrar implicaciones}{Anexo · Semana 2}

\begin{resumen}
Demostrar una afirmación de la forma $P\Longrightarrow Q$ significa justificar que, siempre que se cumple la hipótesis $P$, también se cumple la conclusión $Q$. Una \textbf{demostración directa} parte de $P$ y llega a $Q$ mediante pasos justificados.
\end{resumen}

\begin{nota}[: guía para escribir una demostración directa]
\begin{enumerate}
  \item \textbf{Identifique los objetos y su dominio.} Si la afirmación se refiere a cualesquiera $x_1,x_2\in A$, tome dos elementos arbitrarios de $A$; no elija valores particulares.
  \item \textbf{Distinga hipótesis y conclusión.} Escriba qué afirma $P$ y qué debe obtener para probar $Q$.
  \item \textbf{Suponga $P$.} Esta es la información de partida; todavía no puede dar por cierta la conclusión.
  \item \textbf{Use definiciones y operaciones válidas.} Explique cómo cada paso se obtiene del anterior. Para dividir, compruebe que el divisor no es cero.
  \item \textbf{Concluya $Q$.} Indique qué ha demostrado y por qué el argumento vale para todos los elementos considerados.
\end{enumerate}
\end{nota}

\begin{ejemplo}[Probar inyectividad de una función]
Demuestre que $f:\mathbb{R}\to\mathbb{R}$, $f(x)=3x+1$, es inyectiva. Por definición, debemos probar que, para cualesquiera $x_1,x_2\in\mathbb{R}$,
\[
\underbrace{f(x_1)=f(x_2)}_{P:\ \text{hipótesis}}
\quad\Longrightarrow\quad
\underbrace{x_1=x_2}_{Q:\ \text{conclusión}}.
\]
\end{ejemplo}
\ifmostrarSoluciones
\begin{solucion}
Sean $x_1,x_2\in\mathbb{R}$ arbitrarios y supongamos que $f(x_1)=f(x_2)$. En cada fila, la indicación de la derecha explica cómo pasar a la siguiente igualdad:
\[
\renewcommand{\arraystretch}{1.5}
\begin{array}{rcl@{\qquad}r}
f(x_1) & = & f(x_2)
  & \text{Usamos la ley de asignación de }f. \\
3x_1+1 & = & 3x_2+1
  & \text{Restamos }1\text{ a ambos lados.} \\
3x_1 & = & 3x_2
  & \text{Dividimos entre }3\neq0. \\
x_1 & = & x_2
  & \text{Obtenemos la conclusión }Q.
\end{array}
\]

Así, la hipótesis $P$ implica la conclusión $Q$. Como $x_1$ y $x_2$ eran arbitrarios, la implicación vale para todos los pares del dominio; por tanto, $f$ es inyectiva.
\end{solucion}
\else\vspace{28mm}\fi

\begin{nota}[errores que conviene evitar]
Comprobar algunos pares de números no demuestra una afirmación para todos los pares. Tampoco basta partir de $x_1=x_2$ y obtener $f(x_1)=f(x_2)$: eso prueba la implicación en sentido contrario, que toda función satisface. Para demostrar inyectividad, se parte de la igualdad de imágenes y se deduce la igualdad de entradas.
\end{nota}

\begin{ejercicios}[práctica de implicaciones]
En cada ejercicio, identifique $P$ y $Q$, escriba la hipótesis y justifique los pasos hasta la conclusión.
\begin{enumerate}
  \item Demuestre que $g:\mathbb{R}\to\mathbb{R}$, $g(x)=5x-2$, es inyectiva.
  \item Sean $a,b\in\mathbb{R}$ con $a\neq0$. Demuestre que $h:\mathbb{R}\to\mathbb{R}$, $h(x)=ax+b$, es inyectiva. ¿Dónde utiliza $a\neq0$?
  \item Para $x_1,x_2\in\mathbb{R}$, demuestre la implicación $x_1<x_2\Longrightarrow3x_1+1<3x_2+1$.
\end{enumerate}
\end{ejercicios}
\ifmostrarSoluciones
\begin{solucion}[de los ejercicios]
\begin{enumerate}
  \item Sean $x_1,x_2\in\mathbb{R}$. Si $g(x_1)=g(x_2)$, entonces $5x_1-2=5x_2-2$. Sumando $2$ y dividiendo entre $5$, obtenemos $x_1=x_2$; luego $g$ es inyectiva.
  \item Sean $x_1,x_2\in\mathbb{R}$. Si $h(x_1)=h(x_2)$, entonces $ax_1+b=ax_2+b$. Restamos $b$ y dividimos entre $a$, operación válida porque $a\neq0$. Resulta $x_1=x_2$; luego $h$ es inyectiva.
  \item Sean $x_1,x_2\in\mathbb{R}$ y supongamos $x_1<x_2$. Multiplicar por $3>0$ conserva la desigualdad: $3x_1<3x_2$. Sumar $1$ en ambos miembros da $3x_1+1<3x_2+1$, como se quería demostrar.
\end{enumerate}
\end{solucion}
\fi


\end{document}
